\begin{explanationbox}
	\textbf{\textcolor{CExplanation}{一句话直觉}}

	将三角形边与角的余弦建立等式，通过边算角。

	\medskip
	\textbf{\textcolor{CExplanation}{核心拆解}}
	\begin{description}[style=nextline, leftmargin=2.8em, labelsep=0.5em, itemsep=0.45em]
		\item[\textbf{要点一（公式轮换）：}] 三个公式结构相同，只是边角轮换，记一个即可推出另外两个。
		\item[\textbf{要点二（三边求角）：}] 已知三条边长，代入公式直接计算任意内角的余弦值。
		\item[\textbf{要点三（判形标准）：}] 计算最大角的余弦值，根据符号判定三角形形状（正锐、零直、负钝）。
	\end{description}

	\medskip
	\textbf{\textcolor{CExplanation}{几何本质（边角互化）}}

	余弦定理本身是勾股定理的推广，变形后直接用三边长度表达夹角余弦，使几何角度关系转为纯代数运算。

	\medskip
	\textbf{\textcolor{CExplanation}{代数意义（消角运算）}}

	将角度计算转化为边长平方和与差的比例形式，便于在方程中消去角度，实现边与角的统一。

	\medskip
	\textbf{\textcolor{CExplanation}{考点价值（高频应用）}}
	\begin{itemize}[leftmargin=*, itemsep=0.5em, topsep=0.4em]
		\item \textbf{考法一（直接求角）：} 给出三边数值，求某内角的余弦值和角度大小。
		\item \textbf{考法二（形状判断）：} 给出三边关系或比例，判断三角形锐角/直角/钝角。
	\end{itemize}

	\medskip
	\textbf{\textcolor{CExplanation}{顿悟点}}

	最长边所对的角是最大角，其余弦符号直接决定三角形整体类型；三边不等时先比较大小找到最大角。

	\medskip
	\textbf{\textcolor{CExplanation}{使用场景}}
	\begin{itemize}[leftmargin=*, itemsep=0.5em, topsep=0.4em]
		\item \textbf{场景一（纯边求角）：} 已知三边但未知任何角度，需要求内角余弦或角度值。
		\item \textbf{场景二（形状判定）：} 已知三边长度或比例，需要判断三角形是锐角、直角还是钝角。
	\end{itemize}
\end{explanationbox}
